\documentclass{article}
\title{Holding a source effect until a remote obligation is terminal}
\begin{document}
\maketitle

\paragraph{What changes.}
The previous source-step lease avoids a synchronous remote read only under
a justified bound on source action opportunities. This follow-up replaces
that timing premise with a source promise, a local effect guard, and
positive evidence of remote closure. The physical promise is absence of
the named carrier from the bay, not that the bay is empty of all objects.
The carrier has left and faces back into the bay. Its old command source
asks to return at every opportunity. The guard runs before each native
source step and replaces only a command that would enter the protected
bay. An unrelated carrier's move and radio delivery still run at the first
cut. There is no concurrent source actuator outside the wrapper, no forced
physical displacement and no need for synchronized clocks. A command
already issued to a different actuator would require a separate
effect-side interlock.

\paragraph{Two local interfaces.}
Install a source hold before publishing the capability for the named
operation. The receiving application can spend it while the operation is
live. A successful native callback records paid closure; a failed attempt
can still consume native transmitter budget and does not close the
obligation. An authorized cancellation, when applied at the receiver,
records a fence before making a terminal receipt available; it does not
claim the report was delivered. For this fixture only, either terminal
kind permits return. Arrival of a receipt for this named operation
clears the physical hold. Merely requesting cancellation, or polling
before the fence is installed, does not. After a terminal result the receiver
must continue refusing old uses. No timer can substitute for both obligations
under arbitrary receipt loss. The wrapper has one live capability identity
and models restarts of the application control state, not the underlying
warehouse or broker. A deployed version needs non-reused identities or
an independently justified supersession rule; a transport envelope must
bind the receipt to the right operation and source hold.

The two critical local orders are distinct. At the source, first install
or recover the guard, then accept relevant movement. At the receiver,
an effect and its terminal update occupy a single modeled operation,
while cancellation installs its terminal state before any receipt. An
interruption between effect and terminal recording is not covered by the
atomic modeled step. An actual implementation needs an idempotent effect
or an effect-side commit/recovery record; stable storage of a Boolean alone
does not establish that property.

\paragraph{Invariant and obstruction.}
A nonterminal receiver implies an active source hold. The transition
proof checks this implication from initial issuance through all finite
histories of effect, fence, receipt and other steps, interpreting protected
restart as state preservation. A guarded native return is disallowed when
the hold is active, so every possible admitted remote effect has the
required source clearance. Retiring the guard on cancellation request
violates the implication. Forgetting the guard at source restart, or
forgetting the terminal fence after the source has released, violates it
in opposite directions. Without terminal evidence, a closed receiver and
one still able to use the capability can have the same local source view.
Unconditionally returning on that view cannot guarantee safety. Delaying
only the conflicting movement carries an explicit liveness cost instead:
a lost terminal receipt can leave return blocked indefinitely.

\paragraph{Replay.}
The schedule enumerates orders of offer arrival, two use attempts, request
and application of cancellation, a terminal read opportunity, and a
control restart, respecting request before applied fence and attempt 1
before attempt 2. First-attempt RF loss and three restart scopes are swept.
The existing broker fixture executes active transmitter/receiver checks,
RF behavior, and actual delivery callbacks; the pinned warehouse executes
each admitted or guarded movement. The effects and custody protocol are
not fabricated as a table of returned outcomes. The wrapper still supplies
which messages have arrived and its durable-state semantics.

In 7,560 configurations per arm, the scoped guard gives 2,394 deliveries,
no unauthorized deliveries or duplicates, and all independent first-cut
continuations. It ties an independent same-information route check using
the same pledge. A blanket halt preserves the checked safety and deliveries
but misses every unrelated first-cut continuation. Releasing on cancellation
request permits 1,638 unauthorized deliveries. Forgetting the source guard
on restart permits 882; forgetting terminal receiver state permits 120 and
8 extra callbacks across the sweep (counts can overlap). These are schedule
counts in a finite fixture, not rates. An actual timeout-based best
established adapter is not compared. This establishes a reusable local
contract shape and an explicit availability limit, not acceptance of the
inventory, an independent mission result, or transport to physical robots.
\end{document}
